\section{Questions}
\label{sec:faq}

\paragraph{Why does WFDU use so many cores?} Worst fit puts every
sub-chain on the emptiest core, so it uses every core of the platform. Its
core count measures spreading, not need: use it on hardware that is
already fixed (use case~4), not to size a platform.

\paragraph{A task-chain is late although all its tasks are schedulable.}
Look at its path (\opt{cores}): every communication between cores costs
LET delays, and every inter-ECU one also a network transfer. A chain
whose sensor and actuator are in different zones crosses the network at
least once, whatever the mapping. Try CADU, a larger \opt{sigma}, or check
that the deadline is not below the sum of the periods the LET semantics
imposes.

\paragraph{The mapping is infeasible.} Read the line below the summary
table (\cref{sec:infeasible}): it names the sub-chain and the reason
(memory, core types, largest number of cores, fixed hardware, or a
sub-chain larger than a core).

\paragraph{\opt{frer} is refused.} Frame replication needs two disjoint
routes: use a ring or a custom topology with a second path.

\paragraph{Can I use my own sensor-to-ECU wiring with generated
applications?} Generate the file (\opt{generate}), then edit the
\key{ecu} of the tasks; or write the application yourself.

\paragraph{Are the results exact?} The bounds are worst cases: a
schedulable core never misses a deadline, a network bound is never
exceeded, a chain on time is always on time, under the models of
\cref{sec:concepts} (fixed priorities, LET, the TSN shapers, round-robin
memory). The simulators check this (\cref{sec:validation}). The number of
cores is the result of a heuristic: within a few percent of the optimum
on the applications of the paper (\cref{sec:ex-exact}).

\paragraph{Where do the generated applications come from?} From the
statistics of the WATERS 2015 engine-management benchmark (periods, shares
of the runnables, activation patterns); the appendix of the paper gives
every rule. \opt{seed} selects one.
